\documentclass[12pt,a4paper,oneside]{article} %стиль оформления - статья с параметрами: шрифт 12, формат бумаги А4, печать с одинаковыпи полями на всех страницах - oneside
\usepackage[T2A]{fontenc} %поддержка кириллицы
\usepackage[utf8]{inputenc} %кодировка символов
\usepackage[english,russian]{babel} %пакет поддержки русского языка с элементами английского
\frenchspacing
\usepackage[dvips]{graphicx} %разрешить включение PostScript-графики
\graphicspath{{images/}}
\usepackage{setspace} %"умное" расстояние между строк
\usepackage[usenames,dvipsnames]{color}
\usepackage{hhline} %работа с таблицами
\usepackage[12pt]{extsizes} %размер шрифта
\usepackage{fancyhdr} %установка колонтитулов
%\usepackage{graphics}

\usepackage{geometry} %ручная установка полей
\geometry{top=1cm} %поле сверху
\geometry{bottom=1cm} %поле снизу
\geometry{left=2cm} %поле слева
\geometry{right=2cm} %поле справа

\begin{document}
\singlespacing % межстрочный одиночный интервал для всего текста
\pagestyle{empty} % нумерация выкл.

%Н:задание на уир и кп
\newpage
\begin{center}
\fontfamily{ptm} \fontsize{12.5pt}{15pt}
\selectfont {\renewcommand{\baselinestretch}{1.2}МИНИСТЕРСТВО ОБРАЗОВАНИЯ И НАУКИ РОССИЙСКОЙ ФЕДЕРАЦИИ}

{\fontfamily{ptm} \fontsize{13pt}{24pt}
\selectfont
Федеральное государственное бюджетное образовательное учреждение \linebreak
высшего профессионального образования}
\end{center}

\begin{center}
{\fontfamily{ptm} \fontsize{15pt}{5pt} 
\selectfont \fontseries{b}
"Национальный исследовательский ядерный университет "МИФИ"
\\[10pt]}

\begin{figure}[h]
\includegraphics[width=0.25\linewidth]{mifi}
\end{figure}
{\fontfamily{ptm} \fontsize{12pt}{14pt}
\selectfont \fontseries{b}
ФАКУЛЬТЕТ КИБЕРНЕТИКИ
\\[15pt]
КАФЕДРА ИНФОРМАЦИОННЫЕ ТЕХНОЛОГИИ
\\[50pt]}

{\fontfamily{ptm} \fontsize{24pt}{24pt}
\selectfont \fontseries{b}
Задание на УИР и КП
\\[20pt]}

\begin{flushleft}
{\fontfamily{ptm} \fontsize{14pt}{24pt}
\selectfont
Студентке гр. \underline{К7-361 Лаврентьевой Марине Олеговне\hspace{5cm}}
\\[40pt]}
\end{flushleft}

{\fontfamily{ptm} \fontsize{16pt}{24pt}
\selectfont 
ТЕМА УИР и КП}

\begin{flushleft}
{\fontfamily{ptm} \fontsize{14pt}{20pt}
\selectfont
\underline{Разработка модуля преобразования геоданных из формата\hspace{4cm}} \linebreak
\underline{OpenStreetMap в формат OpenGIS\hspace{9,4cm}}
\\[60pt]}
\end{flushleft}

{\fontfamily{ptm} \fontsize{14pt}{20pt}
\selectfont
ЗАДАНИЕ}
\end{center}

{\fontfamily{ptm} \fontsize{14pt}{20pt}
\selectfont
\begin{enumerate}
\item Изучение структуры xml-документа с геоданными в формате \linebreak
OpenStreetMap.
\item Построение XPath-запросов для извлечения геоданных об объектах \linebreak
заданного типа.
\item Изучение стандарта OpenGIS и разработка структуры базы \linebreak
данных для хранения геоданных в формате OpenGIS.
\item Построение SQL-запросов для сохранения и извлечения геоданных \linebreak
в формате OpenGIS.
\item Реализация модуля преобразования геоданных из формата \linebreak
OpenStreetMap в формат OpenGIS.
\end{enumerate}}

\newpage
\begin{flushleft}
{\fontfamily{ptm} \fontsize{12pt}{20pt}
\selectfont
Место выполнения УИР и КП \underline{\hspace{11cm}}

Руководитель \underline{\hspace{10,9cm} (Муравьёв С.К.)}
\\[10pt]
Дата выдачи задания \underline{\hspace{12,8cm}}}
\end{flushleft}

\newpage
\subsubsection*{ОТЗЫВ О РАБОТЕ СТУДЕНТА}

\vfill
Руководитель \underline{\hspace{8cm}\glqq{\hspace{1cm}}\grqq\hspace{4cm}2011г.}
%К: задание на уир и кп
\end{document}
